\ifdefined\inMainDoc\else
\documentclass[12pt,a4paper]{article}
% ================================================================
%  PRÉAMBULE COMMUN - ANNALES CORRIGÉES
%  Université de Kara - FaST - Licence Fondamentale MPC
%  Design optimisé impression noir et blanc
% ================================================================

% -- ENCODAGE & LANGUE --
\usepackage[utf8]{inputenc}
\usepackage[T1]{fontenc}
\usepackage[french]{babel}
\usepackage{lmodern}
\usepackage{microtype}

% -- MISE EN PAGE --
\usepackage[a4paper, top=2.8cm, bottom=2.5cm, left=2.5cm, right=2.5cm, headheight=14pt]{geometry}
\usepackage{setspace}
\setstretch{1.2}
\usepackage{parskip}
\setlength{\parindent}{1.2em}
\setlength{\parskip}{4pt}

% -- MATHS --
\usepackage{amsmath, amssymb, mathtools, amsthm}
\usepackage{mathrsfs}
\usepackage{dsfont}
\usepackage{bm}
\usepackage{cancel}
\usepackage{textcomp}
% Définition de \oiint (intégrale de surface fermée) sans package supplémentaire
\newcommand{\oiint}{\mathop{\iint\!\!\!\!\!\!\!\bigcirc}\limits}

% -- PHYSIQUE & CHIMIE --
\usepackage{physics}
\usepackage{siunitx}
% Lorsque physics est chargé, siunitx omet \qty. On le restaure ici :
\AtBeginDocument{\RenewCommandCopy\qty\SI}
\sisetup{locale = FR, detect-all, output-decimal-marker={,}}
% Commande de substitution pour les unités posant problème avec pdflatex
% (ohm, degreeCelsius, micro, angstrom, percent utilisent des caractères non-ASCII)
\newcommand{\qtyohm}[1]{\ensuremath{\num{#1}\,\Omega}}
\newcommand{\qtydegc}[1]{\ensuremath{\num{#1}\,{}^\circ\text{C}}}
\newcommand{\qtyangstrom}[1]{\ensuremath{\num{#1}\,\text{\AA}}}
\newcommand{\qtypercent}[1]{\ensuremath{\num{#1}\,\%}}
\newcommand{\qtyum}[1]{\ensuremath{\num{#1}\,\mu\text{m}}}
\usepackage[version=4]{mhchem}

% -- GRAPHISMES --
\usepackage{graphicx}
\usepackage{float}
\usepackage{caption}
\usepackage{subcaption}
\usepackage{wrapfig}
\usepackage{tikz}
\usetikzlibrary{calc, arrows.meta, patterns, shapes.geometric}
\usepackage{pgfplots}
\pgfplotsset{compat=1.18}

% -- TABLEAUX --
\usepackage{booktabs}
\usepackage{array}
\usepackage{tabularx}
\usepackage{multirow}

% -- LISTES --
\usepackage{enumitem}
\setlist[enumerate]{topsep=3pt, itemsep=2pt, parsep=0pt}
\setlist[itemize]{topsep=3pt, itemsep=2pt, parsep=0pt, label={.}}

% -- BOÎTES (noir et blanc) --
\usepackage{mdframed}
\usepackage{tcolorbox}
\tcbuselibrary{skins, breakable, theorems}

% Boîte EXERCICE : encadré avec filet épais, fond blanc
\newmdenv[
  linewidth=1.2pt,
  linecolor=black,
  roundcorner=3pt,
  skipabove=10pt,
  skipbelow=6pt,
  innertopmargin=8pt,
  innerbottommargin=8pt,
  innerleftmargin=10pt,
  innerrightmargin=10pt,
  backgroundcolor=white,
  frametitle={\bfseries\small Exercice},
  frametitlealignment=\raggedright,
  frametitleaboveskip=4pt,
  frametitlebelowskip=4pt,
]{boiteexercice}

% Boîte CORRIGÉ : fond gris très clair + filet fin
\newmdenv[
  linewidth=0.6pt,
  linecolor=black,
  leftline=true,
  rightline=false,
  topline=false,
  bottomline=false,
  linecolor=black,
  innerleftmargin=12pt,
  innerrightmargin=8pt,
  innertopmargin=6pt,
  innerbottommargin=6pt,
  backgroundcolor=gray!8,
  skipabove=4pt,
  skipbelow=8pt,
]{boitecorrige}

% Boîte RÉSUMÉ DE COURS : double encadré
\newmdenv[
  linewidth=0.8pt,
  linecolor=black,
  roundcorner=2pt,
  backgroundcolor=gray!12,
  innertopmargin=10pt,
  innerbottommargin=10pt,
  innerleftmargin=12pt,
  innerrightmargin=12pt,
  skipabove=12pt,
  skipbelow=8pt,
]{boiteresume}

% Boîte ATTENTION/REMARQUE : barre gauche épaisse
\newmdenv[
  leftline=true,
  rightline=false,
  topline=false,
  bottomline=false,
  linewidth=3pt,
  linecolor=black,
  backgroundcolor=gray!5,
  innerleftmargin=12pt,
  innerrightmargin=8pt,
  innertopmargin=5pt,
  innerbottommargin=5pt,
  skipabove=6pt,
  skipbelow=6pt,
]{boiteremarque}

% -- DIFFICULTÉ DES EXERCICES --
\newcommand{\facile}{$\star$}
\newcommand{\moyen}{$\star\!\star$}
\newcommand{\difficile}{$\star\!\star\!\star$}
\newcommand{\niveau}[1]{\hfill\textbf{#1}\hspace{2pt}}

% -- EN-TÊTES ET PIEDS DE PAGE --
\usepackage{fancyhdr}
\pagestyle{fancy}
\fancyhf{}
\renewcommand{\headrulewidth}{0.5pt}
\renewcommand{\footrulewidth}{0.3pt}

% -- HYPERLIENS --
\usepackage[hidelinks]{hyperref}
\hypersetup{
    colorlinks=false,
    pdftitle={Annale Corrigée - Licence MPC - Université de Kara}
}

% -- TITRES --
\usepackage{titlesec}
\titleformat{\section}{\Large\bfseries}{\thesection.}{0.8em}{}[\vspace{2pt}\hrule\vspace{4pt}]
\titleformat{\subsection}{\large\bfseries}{\thesubsection.}{0.7em}{}
\titleformat{\subsubsection}{\normalsize\bfseries}{\thesubsubsection.}{0.6em}{}

% -- THÉORÈMES --
\theoremstyle{definition}
\newtheorem{defn}{Définition}[section]
\newtheorem{prop}{Propriété}[section]
\newtheorem{thm}{Théorème}[section]
\newtheorem{corol}{Corollaire}[section]
\newtheorem{exemple}{Exemple}[section]

% -- COMMANDES MATHÉMATIQUES UTILES --
\newcommand{\vect}[1]{\overrightarrow{#1}}
\newcommand{\R}{\mathbb{R}}
\newcommand{\N}{\mathbb{N}}
\newcommand{\Z}{\mathbb{Z}}
\newcommand{\C}{\mathbb{C}}
\renewcommand{\d}{\mathrm{d}}
\newcommand{\deriv}[2]{\dfrac{\d #1}{\d #2}}
\newcommand{\dpart}[2]{\dfrac{\partial #1}{\partial #2}}
\newcommand{\rot}{\overrightarrow{\mathrm{rot}}\,}
\renewcommand{\div}{\mathrm{div}\,}
\renewcommand{\grad}{\overrightarrow{\mathrm{grad}}\,}
\newcommand{\e}{\mathrm{e}}
\newcommand{\ii}{\mathrm{i}}
\renewcommand{\Tr}{\mathrm{Tr}}

% Numérotation des équations par section
\numberwithin{equation}{section}

% -- RÉSULTATS ENCADRÉS --
\newcommand{\res}[1]{\[\boxed{#1}\]}

% -- REMARQUE INLINE --
\newcommand{\rmq}[1]{\begin{boiteremarque}\textbf{Remarque.} #1\end{boiteremarque}}

% -- MISE EN VALEUR DU RÉSULTAT FINAL --
\newcommand{\reponse}[1]{%
  \begin{center}
    \fbox{\begin{minipage}{0.92\textwidth}\centering #1\end{minipage}}
  \end{center}%
}


\lhead{\small\textit{Programmation Python - Licence 2}}
\rhead{\small\textit{FaST, Université de Kara}}
\cfoot{\thepage}

\usepackage{listings}
\lstset{
  language=Python,
  basicstyle=\ttfamily\small,
  keywordstyle=\bfseries,
  commentstyle=\itshape,
  stringstyle=\ttfamily,
  numbers=left,
  numberstyle=\tiny,
  numbersep=8pt,
  frame=single,
  breaklines=true,
  tabsize=4,
  showstringspaces=false,
  literate={é}{{\'{e}}}1 {è}{{\`{e}}}1 {ê}{{\^{e}}}1
           {à}{{\`a}}1 {ù}{{\`u}}1 {î}{{\^i}}1 {ô}{{\^o}}1
           {ç}{{\c{c}}}1
}

\begin{document}
\fi

\begin{center}
  {\LARGE\bfseries PROGRAMMATION EN PYTHON}\\[0.3cm]
  {\normalsize Licence 2 - Mathématiques, Physique et Chimie}\\[0.1cm]
  \rule{0.7\textwidth}{0.8pt}
\end{center}

\section*{Résumé du cours}

\begin{boiteresume}

\textbf{1. Types et variables.} Python est un langage interprété à typage dynamique. Types de base : \texttt{int}, \texttt{float}, \texttt{complex}, \texttt{str}, \texttt{bool}, \texttt{list}, \texttt{tuple}, \texttt{dict}, \texttt{set}.

\textbf{2. Structures de contrôle.}
\begin{lstlisting}
if condition:           # conditionnel
    ...
elif autre:
    ...
else:
    ...

for i in range(n):      # boucle for
    ...

while condition:        # boucle while
    ...
\end{lstlisting}

\textbf{3. Fonctions.}
\begin{lstlisting}
def nom_fonction(param1, param2, defaut=0):
    """Docstring."""
    return resultat
\end{lstlisting}

\textbf{4. Bibliothèques scientifiques.}
\begin{itemize}[nosep]
  \item \texttt{numpy} : tableaux numériques, algèbre linéaire, fonctions mathématiques.
  \item \texttt{matplotlib} : tracé de courbes (\texttt{plt.plot}, \texttt{plt.show}).
  \item \texttt{scipy} : intégration, optimisation, statistiques.
\end{itemize}

\textbf{5. Listes et compréhensions.} \texttt{[f(x) for x in liste if condition]} crée une liste en une ligne. Les tableaux NumPy (\texttt{np.array}) permettent des opérations vectorisées (plus rapides que les boucles).

\end{boiteresume}

\newpage
\section{Bases du langage}

\subsection*{Exercice 1 - Variables et opérations \niveau{\facile}}

\begin{boiteexercice}
Écrire un programme Python qui :
\begin{enumerate}
  \item Demande à l'utilisateur deux réels $a$ et $b$.
  \item Affiche leurs somme, différence, produit et rapport.
  \item Résout l'équation $ax^2 + bx + c = 0$ (en demandant aussi $c$) en affichant toutes les racines complexes.
\end{enumerate}
\end{boiteexercice}

\begin{boitecorrige}
\begin{lstlisting}
import cmath  # pour les racines complexes

# Partie 1 : operations de base
a = float(input("Entrez a : "))
b = float(input("Entrez b : "))

print(f"Somme      : {a} + {b} = {a + b}")
print(f"Difference : {a} - {b} = {a - b}")
print(f"Produit    : {a} * {b} = {a * b}")
if b != 0:
    print(f"Rapport    : {a} / {b} = {a / b:.4f}")
else:
    print("Division par zero impossible.")

# Partie 2 : trinome du second degre
print("\n- Resolution du trinome ax^2 + bx + c -")
c = float(input("Entrez c : "))

if a == 0:
    print("Ce n'est pas un trinome (a = 0).")
else:
    delta = b**2 - 4*a*c
    x1 = (-b + cmath.sqrt(delta)) / (2*a)
    x2 = (-b - cmath.sqrt(delta)) / (2*a)
    if delta > 0:
        print(f"Deux racines reelles : x1 = {x1.real:.4f}, x2 = {x2.real:.4f}")
    elif delta == 0:
        print(f"Racine double : x = {x1.real:.4f}")
    else:
        print(f"Racines complexes : x1 = {x1}, x2 = {x2}")
\end{lstlisting}
\end{boitecorrige}

\section{Calcul scientifique avec NumPy et Matplotlib}

\subsection*{Exercice 2 - Tracé de fonctions \niveau{\facile}}

\begin{boiteexercice}
Tracer sur le même graphique les fonctions $f(x)=\sin(x)$, $g(x)=\cos(x)$ et $h(x)=\sin(2x)$ pour $x\in[0,2\pi]$ avec une légende, un titre et des labels d'axes. Ajouter une grille.
\end{boiteexercice}

\begin{boitecorrige}
\begin{lstlisting}
import numpy as np
import matplotlib.pyplot as plt

x = np.linspace(0, 2*np.pi, 500)

plt.figure(figsize=(9, 5))
plt.plot(x, np.sin(x),   label=r'$\sin(x)$',   linewidth=2)
plt.plot(x, np.cos(x),   label=r'$\cos(x)$',   linewidth=2, linestyle='--')
plt.plot(x, np.sin(2*x), label=r'$\sin(2x)$',  linewidth=2, linestyle=':')

plt.axhline(0, color='black', linewidth=0.8)
plt.xlabel('$x$ (radians)', fontsize=12)
plt.ylabel('Valeur', fontsize=12)
plt.title('Fonctions trigonometriques', fontsize=14)
plt.legend(fontsize=12)
plt.grid(True, alpha=0.4)
plt.xticks([0, np.pi/2, np.pi, 3*np.pi/2, 2*np.pi],
           ['0', r'$\pi/2$', r'$\pi$', r'$3\pi/2$', r'$2\pi$'])
plt.tight_layout()
plt.show()
\end{lstlisting}

\texttt{np.linspace(0, 2*np.pi, 500)} génère 500 points régulièrement espacés entre 0 et $2\pi$. Les opérations NumPy (\texttt{np.sin}, \texttt{np.cos}) s'appliquent vectoriellement à tous les points simultanément.
\end{boitecorrige}

\subsection*{Exercice 3 - Méthodes numériques \niveau{\moyen}}

\begin{boiteexercice}
\begin{enumerate}
  \item Implémenter la méthode de la bissection pour trouver la racine de $f(x)=x^3-x-2$ sur $[1,2]$.
  \item Implémenter la méthode d'Euler pour résoudre $y'=y\sin(x)$, $y(0)=1$ sur $[0,5]$. Comparer avec la solution exacte $y(x)=\e^{1-\cos x}$.
\end{enumerate}
\end{boiteexercice}

\begin{boitecorrige}
\textbf{1. Méthode de la bissection :}
\begin{lstlisting}
def bissection(f, a, b, tol=1e-8, max_iter=100):
    """Trouve la racine de f sur [a,b] par bissection."""
    if f(a) * f(b) > 0:
        raise ValueError("f(a) et f(b) doivent etre de signes opposes.")
    for _ in range(max_iter):
        c = (a + b) / 2
        if abs(f(c)) < tol or (b - a) / 2 < tol:
            return c
        if f(a) * f(c) < 0:
            b = c
        else:
            a = c
    return (a + b) / 2

f = lambda x: x**3 - x - 2
racine = bissection(f, 1, 2)
print(f"Racine : x = {racine:.8f}")  # x ~= 1.52137971
\end{lstlisting}

\textbf{2. Méthode d'Euler :}
\begin{lstlisting}
import numpy as np
import matplotlib.pyplot as plt

# Methode d'Euler
def euler(f, y0, x0, xf, n):
    h = (xf - x0) / n
    x = np.linspace(x0, xf, n+1)
    y = np.zeros(n+1)
    y[0] = y0
    for i in range(n):
        y[i+1] = y[i] + h * f(x[i], y[i])
    return x, y

f = lambda x, y: y * np.sin(x)
x_euler, y_euler = euler(f, 1, 0, 5, 200)

# Solution exacte
x_exact = np.linspace(0, 5, 1000)
y_exact = np.exp(1 - np.cos(x_exact))

plt.plot(x_euler, y_euler, label="Euler (n=200)", linestyle='--')
plt.plot(x_exact, y_exact, label="Solution exacte")
plt.legend(); plt.grid(True); plt.show()
\end{lstlisting}
\end{boitecorrige}

\subsection*{Exercice 4 - Algèbre linéaire avec NumPy \niveau{\moyen}}

\begin{boiteexercice}
Résoudre le système $A\vec{x}=\vec{b}$ avec :
\[
A = \begin{pmatrix}2&1&-1\\-3&-1&2\\-2&1&2\end{pmatrix}, \quad \vec{b} = \begin{pmatrix}8\\-11\\-3\end{pmatrix}
\]
en utilisant \texttt{numpy.linalg.solve}. Vérifier la solution, calculer $\det(A)$ et les valeurs propres de $A$.
\end{boiteexercice}

\begin{boitecorrige}
\begin{lstlisting}
import numpy as np

A = np.array([[2, 1, -1],
              [-3, -1, 2],
              [-2, 1, 2]], dtype=float)

b = np.array([8, -11, -3], dtype=float)

# Resolution du systeme
x = np.linalg.solve(A, b)
print(f"Solution : x = {x}")  # x = [2, 3, -1]

# Verification : A @ x doit etre egale a b
residuel = np.linalg.norm(A @ x - b)
print(f"Residuel ||Ax - b|| = {residuel:.2e}")

# Determinant
det = np.linalg.det(A)
print(f"det(A) = {det:.4f}")

# Valeurs propres
valeurs_propres, vecteurs_propres = np.linalg.eig(A)
print(f"Valeurs propres : {valeurs_propres}")
\end{lstlisting}

La solution est $x_1=2$, $x_2=3$, $x_3=-1$. \texttt{np.linalg.solve} utilise la décomposition LU, ce qui est bien plus efficace que l'inversion de matrice (\texttt{np.linalg.inv}).
\end{boitecorrige}

\subsection*{Exercice 5 - Traitement de données et statistiques \niveau{\difficile}}

\begin{boiteexercice}
On dispose d'une série de mesures de la résistance d'un conducteur à différentes températures :

\begin{center}
\begin{tabular}{c|cccccc}
$T$ (°C) & 20 & 40 & 60 & 80 & 100 & 120\\
\hline
$R$ ($\Omega$) & 10.2 & 11.1 & 12.0 & 12.8 & 13.9 & 14.7
\end{tabular}
\end{center}

\begin{enumerate}
  \item Calculer la moyenne, l'écart-type et le coefficient de variation de $R$.
  \item Effectuer une régression linéaire $R = aT + b$ par la méthode des moindres carrés (\texttt{numpy.polyfit}).
  \item Estimer la résistance à $T = \qtydegc{150}$.
  \item Tracer les données et la droite de régression.
\end{enumerate}
\end{boiteexercice}

\begin{boitecorrige}
\begin{lstlisting}
import numpy as np
import matplotlib.pyplot as plt

T = np.array([20, 40, 60, 80, 100, 120], dtype=float)
R = np.array([10.2, 11.1, 12.0, 12.8, 13.9, 14.7])

# 1. Statistiques descriptives
print(f"Moyenne de R  : {np.mean(R):.3f} Ohm")
print(f"Ecart-type    : {np.std(R, ddof=1):.3f} Ohm")
print(f"Coeff. variation : {100*np.std(R,ddof=1)/np.mean(R):.1f} %")

# 2. Regression lineaire
coeffs = np.polyfit(T, R, 1)
a, b = coeffs
print(f"\nRegression : R = {a:.4f} * T + {b:.4f}")

# Coefficient de correlation
r = np.corrcoef(T, R)[0, 1]
print(f"Coefficient de correlation r = {r:.4f}")

# 3. Estimation a T = 150 degC
T_pred = 150
R_pred = a * T_pred + b
print(f"\nEstimation a {T_pred} degC : R = {R_pred:.2f} Ohm")

# 4. Trace
T_fit = np.linspace(0, 160, 100)
R_fit = a * T_fit + b

plt.scatter(T, R, color='black', zorder=5, label='Mesures')
plt.plot(T_fit, R_fit, label=f'Regression: R={a:.3f}T+{b:.3f}')
plt.xlabel('Temperature (degC)'); plt.ylabel('Resistance (Ohm)')
plt.title('Resistance vs Temperature')
plt.legend(); plt.grid(True); plt.show()
\end{lstlisting}

Résultats : $a \approx \ensuremath{\num{0.045}\,\Omega/{}^\circ\text{C}}$, $b \approx \qty{9.28}{\Omega}$, $r \approx 0{,}9997$. La linéarité est excellente. À 150°C, $R\approx\qty{16.0}{\Omega}$.
\end{boitecorrige}

\section{Structures de données et algorithmes}

\subsection*{Exercice 6 - Listes, tri et compréhension \niveau{\facile}}

\begin{boiteexercice}
\begin{enumerate}
  \item Créer une liste des carrés des entiers de 1 à 20 en utilisant une compréhension de liste.
  \item Filtrer les éléments pairs de cette liste.
  \item Trier une liste de mots par longueur (puis alphabétiquement en cas d'égalité) : \texttt{mots = ["chat", "python", "code", "AI", "Numpy", "ax"]}.
  \item Écrire une fonction qui recherche un élément dans une liste triée par l'algorithme de \emph{recherche dichotomique}.
\end{enumerate}
\end{boiteexercice}

\begin{boitecorrige}
\begin{lstlisting}
# 1. Comprehension de liste : carres
carres = [n**2 for n in range(1, 21)]
print("Carres:", carres)

# 2. Filtrage des elements pairs
pairs = [x for x in carres if x % 2 == 0]
print("Carres pairs:", pairs)

# 3. Tri par longueur puis alphabetique
mots = ["chat", "python", "code", "AI", "Numpy", "ax"]
mots_tries = sorted(mots, key=lambda m: (len(m), m.lower()))
print("Mots tries:", mots_tries)
# Resultat: ['AI', 'ax', 'chat', 'code', 'Numpy', 'python']

# 4. Recherche dichotomique
def recherche_dicho(liste, cible):
    """Recherche dichotomique dans une liste triee.
    Retourne l'indice si trouve, -1 sinon."""
    gauche, droite = 0, len(liste) - 1
    while gauche <= droite:
        milieu = (gauche + droite) // 2
        if liste[milieu] == cible:
            return milieu
        elif liste[milieu] < cible:
            gauche = milieu + 1
        else:
            droite = milieu - 1
    return -1

nombres = sorted([3, 7, 1, 15, 9, 4, 12])  # [1, 3, 4, 7, 9, 12, 15]
print(f"Recherche de 9 : indice {recherche_dicho(nombres, 9)}")   # 4
print(f"Recherche de 6 : indice {recherche_dicho(nombres, 6)}")   # -1
\end{lstlisting}

\textbf{Complexité :} La recherche dichotomique a une complexité $O(\log_2 n)$ au lieu de $O(n)$ pour la recherche linéaire. Pour $n = 10^6$ éléments, seulement $\sim 20$ comparaisons suffisent.
\end{boitecorrige}

\subsection*{Exercice 7 - Fonctions récursives \niveau{\moyen}}

\begin{boiteexercice}
\begin{enumerate}
  \item Écrire une fonction récursive calculant $n!$ et une version itérative. Comparer pour $n=10$.
  \item Écrire une fonction récursive pour les nombres de Fibonacci $F_n = F_{n-1}+F_{n-2}$, $F_0=0$, $F_1=1$.
  \item Pourquoi la version naïve de Fibonacci est-elle exponentiellement lente ? Implémenter une version mémoïsée avec un dictionnaire.
  \item Calculer $F_{35}$ avec les deux versions et mesurer le temps avec \texttt{time.time()}.
\end{enumerate}
\end{boiteexercice}

\begin{boitecorrige}
\begin{lstlisting}
import time

# 1. Factorielle
def factorielle_rec(n):
    """Factorielle recursive."""
    if n <= 1:
        return 1
    return n * factorielle_rec(n - 1)

def factorielle_iter(n):
    """Factorielle iterative."""
    result = 1
    for i in range(2, n + 1):
        result *= i
    return result

print(f"10! (recursif)  = {factorielle_rec(10)}")
print(f"10! (iteratif)  = {factorielle_iter(10)}")

# 2. Fibonacci naif (exponentiel)
def fib_naif(n):
    if n <= 1:
        return n
    return fib_naif(n-1) + fib_naif(n-2)

# 3. Fibonacci memoize (lineaire)
def fib_memo(n, cache={}):
    if n in cache:
        return cache[n]
    if n <= 1:
        return n
    cache[n] = fib_memo(n-1, cache) + fib_memo(n-2, cache)
    return cache[n]

# 4. Comparaison de temps
n = 35
t0 = time.time()
res_naif = fib_naif(n)
t1 = time.time()
print(f"F({n}) naif    = {res_naif} en {t1-t0:.3f} s")

t0 = time.time()
res_memo = fib_memo(n)
t1 = time.time()
print(f"F({n}) memoize = {res_memo} en {t1-t0:.6f} s")
\end{lstlisting}

\textbf{Analyse :} La version naïve calcule $\sim 2^{35}$ appels (exponentiel), tandis que la version mémoïsée fait $O(n)$ calculs. Pour $n=35$, la version naïve prend $\sim$ 5 secondes, la mémoïsée $\sim$ 0 ms.
\end{boitecorrige}

\subsection*{Exercice 8 - Lecture et écriture de fichiers \niveau{\facile}}

\begin{boiteexercice}
On dispose d'un fichier texte \texttt{mesures.txt} contenant des lignes de la forme \texttt{t R} (temps et résistance).
\begin{enumerate}
  \item Écrire un programme qui génère ce fichier avec des données simulées ($T = [20, 40, 60, 80, 100]$°C, $R = 10 + 0.05T + \text{bruit}$).
  \item Lire le fichier et extraire les données dans deux listes.
  \item Calculer la moyenne et l'écart-type de $R$ à partir des données lues.
  \item Réécrire un fichier \texttt{resultats.txt} avec les statistiques calculées.
\end{enumerate}
\end{boiteexercice}

\begin{boitecorrige}
\begin{lstlisting}
import numpy as np

# 1. Generation et ecriture du fichier
np.random.seed(42)
T_vals = [20, 40, 60, 80, 100]
R_vals = [10 + 0.05*T + np.random.normal(0, 0.1) for T in T_vals]

with open("mesures.txt", "w") as f:
    f.write("# Temperature(C)  Resistance(Ohm)\n")
    for T, R in zip(T_vals, R_vals):
        f.write(f"{T:.1f}  {R:.4f}\n")

print("Fichier mesures.txt cree.")

# 2. Lecture du fichier
temperatures = []
resistances = []

with open("mesures.txt", "r") as f:
    for ligne in f:
        if ligne.startswith("#"):  # ignorer commentaires
            continue
        parts = ligne.strip().split()
        if len(parts) == 2:
            temperatures.append(float(parts[0]))
            resistances.append(float(parts[1]))

print("Temperatures lues:", temperatures)
print("Resistances lues :", [f"{r:.4f}" for r in resistances])

# 3. Statistiques
R_arr = np.array(resistances)
moyenne = np.mean(R_arr)
ecart_type = np.std(R_arr, ddof=1)
print(f"\nMoyenne R = {moyenne:.4f} Ohm")
print(f"Ecart-type = {ecart_type:.4f} Ohm")

# 4. Ecriture des resultats
with open("resultats.txt", "w") as f:
    f.write("=== Resultats statistiques ===\n")
    f.write(f"Nombre de mesures : {len(resistances)}\n")
    f.write(f"Moyenne R         : {moyenne:.4f} Ohm\n")
    f.write(f"Ecart-type        : {ecart_type:.4f} Ohm\n")
    f.write(f"Min / Max         : {R_arr.min():.4f} / {R_arr.max():.4f} Ohm\n")

print("\nFichier resultats.txt ecrit.")
\end{lstlisting}
\end{boitecorrige}

\subsection*{Exercice 9 - NumPy : tableaux et algèbre linéaire \niveau{\moyen}}

\begin{boiteexercice}
\begin{enumerate}
  \item Créer un tableau NumPy $A$ de taille $3\times3$ représentant la matrice de rotation d'angle $\theta = 45°$.
  \item Calculer $\det(A)$, $A^{-1}$, et les valeurs propres et vecteurs propres de $A$.
  \item Résoudre le système $Ax = b$ avec $b = (1, 0, 1)^T$.
  \item Vérifier que $A^{-1}A = I$ et que la solution satisfait bien $Ax = b$.
\end{enumerate}
Matrice de rotation 3D autour de $z$ :
\[
A = \begin{pmatrix}\cos\theta & -\sin\theta & 0\\\sin\theta & \cos\theta & 0\\0 & 0 & 1\end{pmatrix}
\]
\end{boiteexercice}

\begin{boitecorrige}
\begin{lstlisting}
import numpy as np

theta = np.pi / 4  # 45 degres

# 1. Matrice de rotation
A = np.array([
    [np.cos(theta), -np.sin(theta), 0],
    [np.sin(theta),  np.cos(theta), 0],
    [0,              0,             1]
])
print("Matrice A :")
print(np.round(A, 6))

# 2. Determinant, inverse, valeurs propres
det_A = np.linalg.det(A)
print(f"\ndet(A) = {det_A:.6f}")  # doit valoir 1

A_inv = np.linalg.inv(A)
print("\nA^(-1) :")
print(np.round(A_inv, 6))

vals_propres, vects_propres = np.linalg.eig(A)
print(f"\nValeurs propres : {np.round(vals_propres, 4)}")
# Pour une rotation, les VP sont 1 et exp(+-i*theta)

# 3. Resolution du systeme
b = np.array([1.0, 0.0, 1.0])
x = np.linalg.solve(A, b)
print(f"\nSolution x = {np.round(x, 6)}")

# 4. Verification
identite = np.round(A_inv @ A, 10)
print("\nA^(-1) @ A =")
print(identite)

residuel = np.linalg.norm(A @ x - b)
print(f"\nResidu ||Ax - b|| = {residuel:.2e}")
\end{lstlisting}

\textbf{Résultats attendus :} $\det(A) = 1$ (rotation préserve les volumes), valeurs propres $\{1, \e^{i\pi/4}, \e^{-i\pi/4}\}$, résidu $\approx 10^{-15}$.
\end{boitecorrige}

\subsection*{Exercice 10 - Matplotlib : visualisation et histogramme \niveau{\moyen}}

\begin{boiteexercice}
\begin{enumerate}
  \item Générer $N = 1000$ réalisations d'une variable aléatoire gaussienne $X \sim \mathcal{N}(\mu=5, \sigma=2)$ avec NumPy.
  \item Tracer l'histogramme normalisé (densité de probabilité) et superposer la densité théorique $f(x) = \dfrac{1}{\sigma\sqrt{2\pi}}\e^{-(x-\mu)^2/(2\sigma^2)}$.
  \item Calculer la moyenne et l'écart-type empiriques et les comparer aux valeurs théoriques.
  \item Tracer une figure avec deux sous-graphiques : (a) l'histogramme, (b) la courbe de distribution cumulée empirique (ECDF).
\end{enumerate}
\end{boiteexercice}

\begin{boitecorrige}
\begin{lstlisting}
import numpy as np
import matplotlib.pyplot as plt
from scipy.stats import norm

# 1. Generation des donnees
np.random.seed(42)
mu, sigma = 5, 2
N = 1000
data = np.random.normal(mu, sigma, N)

# 3. Statistiques empiriques
print(f"Moyenne empirique : {data.mean():.3f}  (theorique : {mu})")
print(f"Ecart-type empirique : {data.std():.3f}  (theorique : {sigma})")

# 2. et 4. Graphiques
fig, (ax1, ax2) = plt.subplots(1, 2, figsize=(12, 5))

# Sous-graphique 1 : Histogramme
ax1.hist(data, bins=40, density=True, alpha=0.6,
         color='steelblue', label=f'Histogramme (N={N})')

# Densite theorique
x_theo = np.linspace(data.min(), data.max(), 200)
ax1.plot(x_theo, norm.pdf(x_theo, mu, sigma),
         'r-', linewidth=2,
         label=rf'$\mathcal{{N}}(\mu={mu}, \sigma={sigma})$')
ax1.set_xlabel('x'); ax1.set_ylabel('Densite')
ax1.set_title('Histogramme normalise')
ax1.legend()
ax1.grid(alpha=0.3)

# Sous-graphique 2 : ECDF
data_sorted = np.sort(data)
ecdf = np.arange(1, N+1) / N

ax2.step(data_sorted, ecdf, color='steelblue', label='ECDF empirique')
ax2.plot(x_theo, norm.cdf(x_theo, mu, sigma),
         'r-', linewidth=2, label='CDF theorique')
ax2.set_xlabel('x'); ax2.set_ylabel('P(X <= x)')
ax2.set_title('Fonction de repartition')
ax2.legend()
ax2.grid(alpha=0.3)

plt.tight_layout()
plt.savefig("statistiques.pdf", bbox_inches='tight')
plt.show()
\end{lstlisting}

\textbf{Interprétation :} L'histogramme normalisé (densité) converge vers la densité théorique en $1/\sqrt{N}$. L'ECDF (Empirical Cumulative Distribution Function) permet de visualiser la distribution sans choisir une largeur de bin.
\end{boitecorrige}

\section{Algorithmique et structures de données}

\subsection*{Exercice 11 - Parité et année bissextile \niveau{\facile}}

\begin{boiteexercice}
\begin{enumerate}
  \item Écrire un programme qui demande un entier $n$ à l'utilisateur et affiche s'il est pair ou impair.
  \item Écrire un programme qui demande une année et indique si elle est bissextile. Rappel : une année est bissextile si elle est divisible par 4 mais pas par 100, sauf si elle est divisible par 400.
\end{enumerate}
\end{boiteexercice}

\begin{boitecorrige}
\begin{lstlisting}
# 1. Test de parite
n = int(input("Entrez un entier n : "))
if n % 2 == 0:
    print(f"{n} est pair.")
else:
    print(f"{n} est impair.")

# 2. Annee bissextile
annee = int(input("Entrez une annee : "))
if (annee % 4 == 0 and annee % 100 != 0) or (annee % 400 == 0):
    print(f"{annee} est bissextile.")
else:
    print(f"{annee} n'est pas bissextile.")
\end{lstlisting}

\textbf{Tests :} 2024 (bissextile, divisible par 4 mais pas 100), 1900 (non bissextile, divisible par 100 mais pas 400), 2000 (bissextile, divisible par 400).
\end{boitecorrige}

\subsection*{Exercice 12 - PGCD, PPCM et décomposition binaire \niveau{\moyen}}

\begin{boiteexercice}
\begin{enumerate}
  \item Écrire une fonction \texttt{pgcd(a, b)} utilisant l'algorithme d'Euclide, puis une fonction \texttt{ppcm(a, b)} qui en déduit le PPCM via la relation $a \times b = \mathrm{pgcd}(a,b) \times \mathrm{ppcm}(a,b)$.
  \item Écrire une fonction \texttt{binaire(n)} qui retourne la représentation binaire (chaîne de caractères) d'un entier positif sans utiliser \texttt{bin()}.
  \item Tester avec $a = 84$, $b = 60$ et $n = 77$.
\end{enumerate}
\end{boiteexercice}

\begin{boitecorrige}
\begin{lstlisting}
def pgcd(a, b):
    """Algorithme d'Euclide iteratif."""
    while b != 0:
        a, b = b, a % b
    return a

def ppcm(a, b):
    """PPCM via la relation a*b = pgcd(a,b) * ppcm(a,b)."""
    if a == 0 or b == 0:
        return 0
    return abs(a * b) // pgcd(a, b)

def binaire(n):
    """Decomposition binaire d'un entier positif."""
    if n == 0:
        return "0"
    bits = []
    while n > 0:
        bits.append(str(n % 2))
        n //= 2
    return "".join(reversed(bits))

# Tests
a, b = 84, 60
print(f"pgcd({a}, {b}) = {pgcd(a, b)}")   # 12
print(f"ppcm({a}, {b}) = {ppcm(a, b)}")   # 420

n = 77
print(f"binaire({n}) = {binaire(n)}")     # 1001101
print(f"verification : bin({n}) = {bin(n)[2:]}")
\end{lstlisting}

\textbf{Résultats :} $\mathrm{pgcd}(84,60)=12$, $\mathrm{ppcm}(84,60)=420$, $77 = 64+8+4+1 = (1001101)_2$.
\end{boitecorrige}

\subsection*{Exercice 13 - Permutation et équation du premier degré \niveau{\facile}}

\begin{boiteexercice}
\begin{enumerate}
  \item Écrire une fonction \texttt{permute(a, b)} qui échange les valeurs de deux variables sans utiliser de variable temporaire (astuce arithmétique ou \texttt{a, b = b, a}).
  \item Écrire un programme qui résout l'équation $ax + b = 0$ (avec $a, b$ saisis par l'utilisateur) en traitant tous les cas : solution unique, infinité de solutions, aucune solution.
\end{enumerate}
\end{boiteexercice}

\begin{boitecorrige}
\begin{lstlisting}
# 1. Permutation sans variable temporaire
def permute_arith(a, b):
    """Permutation par somme/difference (entiers/reels)."""
    a = a + b
    b = a - b
    a = a - b
    return a, b

def permute_python(a, b):
    """Permutation pythonique (tuples)."""
    return b, a

x, y = 5, 12
x, y = permute_python(x, y)
print(f"Apres permutation : x = {x}, y = {y}")  # x=12, y=5

# 2. Equation du premier degre a*x + b = 0
a = float(input("Coeff a : "))
b = float(input("Coeff b : "))

if a == 0:
    if b == 0:
        print("Infinité de solutions (tout réel x convient).")
    else:
        print("Aucune solution (équation impossible).")
else:
    x = -b / a
    print(f"Solution unique : x = {x:.6f}")
\end{lstlisting}

\rmq{La permutation arithmétique \texttt{a=a+b; b=a-b; a=a-b} ne fonctionne qu'avec des nombres et peut provoquer des dépassements de capacité dans certains langages ; l'échange par tuples \texttt{a, b = b, a} est la méthode idiomatique en Python.}
\end{boitecorrige}

\subsection*{Exercice 14 - Indice et valeur du maximum dans une liste \niveau{\facile}}

\begin{boiteexercice}
Écrire un programme qui :
\begin{enumerate}
  \item Demande à l'utilisateur $n$ nombres réels et les stocke dans une liste.
  \item Détermine la valeur maximale et son indice (première occurrence).
  \item Détermine également la valeur minimale et son indice.
  \item Affiche la moyenne, la médiane et l'étendue.
\end{enumerate}
\end{boiteexercice}

\begin{boitecorrige}
\begin{lstlisting}
# Saisie de la liste
n = int(input("Combien de valeurs ? "))
valeurs = [float(input(f"valeur {i+1} : ")) for i in range(n)]

# Recherche du maximum et de son indice
val_max = valeurs[0]
idx_max = 0
for i, v in enumerate(valeurs):
    if v > val_max:
        val_max = v
        idx_max = i

# Recherche du minimum et de son indice
val_min = valeurs[0]
idx_min = 0
for i, v in enumerate(valeurs):
    if v < val_min:
        val_min = v
        idx_min = i

# Statistiques
moyenne = sum(valeurs) / n
etendue = val_max - val_min
valeurs_triees = sorted(valeurs)
if n % 2 == 1:
    mediane = valeurs_triees[n // 2]
else:
    mediane = (valeurs_triees[n//2 - 1] + valeurs_triees[n//2]) / 2

print(f"Maximum : {val_max} a l'indice {idx_max}")
print(f"Minimum : {val_min} a l'indice {idx_min}")
print(f"Moyenne : {moyenne:.4f}")
print(f"Mediane : {mediane}")
print(f"Etendue : {etendue}")
\end{lstlisting}

\textbf{Note :} En Python, on peut aussi utiliser les fonctions natives \texttt{max()}, \texttt{min()}, \texttt{list.index()} et le module \texttt{statistics} (fonction \texttt{median()}). L'implémentation manuelle reste utile pour comprendre la logique algorithmique.
\end{boitecorrige}

\subsection*{Exercice 15 - Somme de série géométrique et intégrale numérique \niveau{\moyen}}

\begin{boiteexercice}
\begin{enumerate}
  \item Calculer la somme partielle $S_n = \displaystyle\sum_{k=0}^{n} q^k$ pour $q = 0{,}5$ et $n = 20$, et la comparer à la limite $S_\infty = \dfrac{1}{1-q}$. Afficher l'erreur relative.
  \item Calculer numériquement l'intégrale $I = \displaystyle\int_0^1 \e^{-x^2}\,\mathrm{d}x$ par la méthode des rectangles (à gauche) avec $N = 1000$ intervals. Comparer avec la valeur de référence $I \approx 0{,}746824$.
  \item Refaire le calcul avec la méthode des trapèzes et comparer la précision.
\end{enumerate}
\end{boiteexercice}

\begin{boitecorrige}
\begin{lstlisting}
import math

# 1. Somme de la serie geometrique
q = 0.5
n = 20
S_n = sum(q**k for k in range(n + 1))
S_inf = 1 / (1 - q)
erreur_rel = abs(S_n - S_inf) / S_inf * 100
print(f"S_{n}      = {S_n:.10f}")
print(f"S_inf     = {S_inf:.10f}")
print(f"Erreur rel. = {erreur_rel:.3e} %")

# 2. Methode des rectangles (gauche)
def f(x):
    return math.exp(-x**2)

N = 1000
a, b = 0.0, 1.0
h = (b - a) / N
I_rect = sum(f(a + i * h) for i in range(N)) * h
print(f"\nIntegrale (rectangles) : I = {I_rect:.6f}")

# 3. Methode des trapezes
I_trap = (f(a) + f(b)) / 2
for i in range(1, N):
    I_trap += f(a + i * h)
I_trap *= h
print(f"Integrale (trapezes)   : I = {I_trap:.6f}")

I_ref = 0.746824
print(f"Valeur de reference    : I = {I_ref:.6f}")
print(f"Erreur rectangles : {abs(I_rect - I_ref):.2e}")
print(f"Erreur trapezes   : {abs(I_trap - I_ref):.2e}")
\end{lstlisting}

\textbf{Résultats :} $S_{20} \approx 0{,}9999990463$ avec une erreur relative de $\sim 5 \times 10^{-6}$ \%. Pour l'intégrale, la méthode des rectangles donne $I \approx 0{,}746980$ (erreur $\sim 10^{-4}$) et la méthode des trapèzes $I \approx 0{,}746824$ (erreur $\sim 10^{-7}$, plus précise car la méthode des trapèzes est d'ordre 2).
\end{boitecorrige}

\ifdefined\inMainDoc\else\end{document}\fi
